%% kapa: exam=1 round=20 year=2009 period=1 form=A subject=civil_law subject_ko=민법 seq=1 number=1
%% source: Q-Net 감정평가사 기출문제 (cst00309) artlSeq=5024178 file=extracted/1008/5024178/2009년 제20회 감정평가사 1차시험 기출문제/1교시.hwp
%% conversion: kordoc 4.24.51 -> tools/gen1.py (auto); review=proofread-sonnet; flags=none
\begin{question}{1}
민법의 법원(法源)에 관한 설명으로 옳지 \uline{않은} 것은?
\begin{choices}
  \item 민법 제1조의 법률은 형식적 의미의 법률만을 의미하는 것은 아니다.
  \item 삼권분립을 강조하면 판례의 법원성을 인정하기 어렵다.
  \item 헌법상 대통령의 긴급재정경제명령은 법원이 될 수 없다.
  \item 지방자치단체의 조례와 규칙도 법원이 될 수 있다.
  \item 판례는 관습법의 보충적 효력을 인정하고 있다.
\end{choices}
\end{question}
\answer{3}
